\documentclass[a4paper,nofonts,nohyper,justified,colorinlistoftodos,hidelinks]{tufte-book}

\usepackage[T1]{fontenc}
\usepackage[utf8]{inputenc}
\usepackage[czech]{babel}

\usepackage[a-2u]{pdfx}

\usepackage{graphicx}
\usepackage{algorithmicx}
\usepackage{algpseudocode}
\usepackage{algorithm}
\usepackage{color}
\usepackage{tikz}
\usepackage{pgfplots}
\usepackage{makeidx}
\usepackage{subfigure}

\usepackage{todonotes}

\usepackage{libertine}
\usepackage{libertinust1math}
\usepackage[scaled=0.83]{beramono}
\usepackage{microtype}
\usepackage{xfrac}
\usepackage{xspace}

\graphicspath{{./graphics/}}
 
\floatname{algorithm}{Algoritmus}
\renewcommand{\listalgorithmname}{Seznam algoritmů}

% mathematical notation
\newcommand{\Real}{{\ensuremath{\mathcal{R}}}}
\newcommand{\normal}[2]{\mathcal{N}(#1, #2)}
\newcommand{\uniform}[2]{\mathcal{U}(#1, #2)}
\newcommand{\twopartdef}[4]
{
	\left\{
		\begin{array}{ll}
			#1 & \text{pro } #2 \\
			#3 & \text{pro } #4
		\end{array}
	\right.
}
\newcommand{\twopartdefotherwise}[4]
{
	\left\{
		\begin{array}{ll}
			#1 & \text{pro } #2 \\
			#3 & \text{jinak} #4
		\end{array}
	\right.
}


\pgfplotsset{vasymptote/.style={
    before end axis/.append code={
        \draw[densely dashed] ({rel axis cs:0,0} -| {axis cs:#1,0})
        -- ({rel axis cs:0,1} -| {axis cs:#1,0});
    }
}}

\newcommand{\de}[3]{\texttt{DE/#1/#2/#3}}

% todos, comments, ...
\newcommand{\pozn}[1]{\todo[inline,color={CornflowerBlue},size=\small]{\sc{#1}}}
\newcommand{\missingref}{\todo[color={WildStrawberry}]{\textbf{Missing ref.}}}
\newcommand{\add}[1]{\todo[inline, color={YellowGreen}]{#1}}
\newcommand{\ask}[1]{\todo[inline, color={Peach}]{\textbf{???} #1}}
\newcommand{\redo}[1]{\todo[inline, color={BurntOrange}]{\textbf{<!>} #1}}

\newcommand{\code}[1]{\texttt{#1}}


\makeatletter
\renewcommand{\maketitlepage}{%
\begingroup%
\setlength{\parindent}{0pt}

\includegraphics[width=\textwidth]{graphics/logotyp_fakulty1}

\vspace{1.25in}{\fontsize{36}{54}\selectfont\@title\par}
\vspace{0.5in}{\fontsize{24}{24}\selectfont\textit{\@author}\par}

\vspace{0.5in}{\fontsize{14}{14}\selectfont\textsf{\smallcaps{\@date}}\par}

\vfill{\fontsize{14}{14}\selectfont\textit{\@publisher}\par}

\thispagestyle{empty}
\endgroup
}
\makeatother

\makeindex

\frontmatter

\title{Přírodou inspirované prohledávací algoritmy}
\author{Roman Neruda, Martin Pilát}
\date{\today}

\begin{document}

\maketitle

\mainmatter

%\chapter{Před úvodem}
%\section{Jak to je?}
%\pozn{osnova prednasky tak jak je ted}

Evoluční algoritmy

Jednoduchý genetický algoritmus

Teorie schémat

Reprezentace a operátory v GA

Evoluce kooperace

Evoluční strategie

Diferenciální evoluce

Particle swarm optimization

Evoluční strojové učení

Vícekriteriální optimalizace

Evoluční kombinatorická optimalizace

Ladění, řízení, metaevoluce

Teorie EVA podruhé

Evoluční programování

Genetické programování

Neuroevoluce

Memetické algoritmy

Dynamické krajiny fitness

Tabu search, scatter search

Biologicky věrnější evoluce


%\section{Jak to bude?}
%\section{Pok pok pokusy}
%\pozn{skoro to tu muzeme smazat ne?}

\index{Genetický algoritmus}

\begin{algorithm}
\caption{Schéma evolučního algoritmu}
\label{obreva2}
\begin{algorithmic}
\Procedure{Evoluční algoritmus}{} 
\State $t \gets 0$
\State \emph{Inicializuj} populaci $P_t$ náhodně vygenerovanými jedinci
\State \emph{Ohodnoť} jedince v populaci $P_t$
\While{neplatí \emph{kritérium ukončení}}
\State 	vyber z $P_t$ rodiče $P'_{t+1}$ \emph{Rodičovskou selekcí}
\State 	\emph{Rekombinací} rodičů vzniknou potomci $P'_{t+1}$
\State 	\emph{Mutuj} potomky $P'_{t+1}$
\State 	\emph{Ohodnoť} potomky $P'_{t+1}$
\State 	\emph{Enviromentální selekcí} vyber $P_{t+1}$ z $P_t$ a $P'_{t+1}$
\State $t \gets t+1$
\EndWhile
%\Until{kritérium ukončení}
\EndProcedure
\end{algorithmic}
\end{algorithm}



\chapter{Úvod}

Prohled�v�n� prostoru �e�en� parametricky zadan�ch �loh je jedn�m z hlavn�ch probl�m� mnoha oblast� informatiky i matematiky. Aplikace prohled�v�n� jsou nes�etn� a zahrnuj� optimalizaci re�ln�ch parametr�, kombinatorickou optimalizaci, automatick� v�voj program�, u�en� robot� nebo n�vrh architektury a u�en� neuronov�ch s�t�. P��rodou inspirovan� algoritmy p�inesly do tohoto oboru n�kolik nov�ch p��stup�, kter� jsou schopny efektivn� pracovat i ve vysoce dimenzion�ln�ch prostorech, optimalizovat neline�rn� funkce s lok�ln�mi extr�my a pracovat v tak strukturovan�ch prostorech jako jsou syntaktick� stromy, grafy nebo neuronov� s�t�. 

V�hodou v�t�iny p��stup� zm�n�n�ch v tomto �l�nku je jejich obecnost. Evolu�n� i rojov� techniky m��eme ch�pat jako metaheuristiky, kter� nevy�aduj� mnoho specifick�ch informac� o �e�en� �loze, krom� ohodnocen� jednotliv�ch kandid�t� �e�en� prost�ednictv�m ��elov� funkce naz�van� t� ohodnocovac� funkce nebo fitness. V�hodou tohoto p��stupu je snadn� aplikovatelnost generick�ho algoritmu na �irok� spektrum �loh. Na druhou stranu, existuj� teoretick� i praktick� d�vody pro p�izp�soben� obecn�ch prohled�vac�ch heuristik dan� �loze. Takto p�izp�soben� metody obvykle dosahuj� lep��ch v�sledk� a z�rove� si uchov�vaj� obecn� v�hody t�chto technik, jako je odolnost proti uv�znut� v lok�ln�ch extr�mech. 

Hlavn�m rysem zde uv�d�n�ch p��rodou inspirovan�ch prohled�vac�ch technik je popula�n� p��stup. Na rozd�l od v�t�iny algoritm� lok�ln�ho prohled�van�, kter� zkoumaj� nejbli��� okol� jednoho bodu v prohled�vac�m prostoru, evolu�n� techniky pracuj� s populac� des�tek a� tis�c� �e�en�, kter� paraleln� prohled�vaj� prostor �e�en� a navz�jem se ovliv�uj�. 

Lok�ln� prohled�vac� metody �it� na m�ru ur�it�mu probl�mu typicky maj� v�hodu rychlej��ho lok�ln�ho prohled�v�n� d�ky vyu�it� specifick�ch informac� o probl�mu, jako je nap��klad informace o gradientu ��elov� funkce. V dal��m textu se zm�n�me tak� o mo�nosti hybridizovat n�kter� evolu�n� algoritmy pomoc� specifick�ho lok�ln�ho prohled�v�n�. Tato technika se v~praxi osv�d�uje zrychlen�m konvergence algoritmu, i kdy� n�kdy p�in�� v�t�� nebezpe�� uv�znut� v lok�ln�ch extr�mech ��elov� funkce. Je zaj�mav�, �e z~hlediska p�vodn� biologick� inspirace p�edstavuje tato hybridizace p�ekro�en� darwinistick�ho r�mce a p�in�� lamarckistick� �i epigenetick� prinicpipy. 

\redo{update konec uvodu}
V n�sleduj�c�m textu nejprve stru�n� zm�n�me inspiraci a obecn� rysy evolu�n�ch algoritm� a pak se budeme v�novat n�kolika konkr�tn�m oblastem vych�zej�c�ch z t�chto obecn�ch princip�. Nejprve pop��eme tradi�n� genetick� algoritmy pracuj�c� p�vodn� nad bin�rn� zak�dovan�mi jedinci. D�le budeme hovo�it o evolu�n�m programov�n�, oblasti kter� st�r� rozd�l mezi genotypem (zak�dov�n�m jedince pro ��ely evolu�n�ho prohled�v�n�) a fenotypem (vlastn�m modelem, kter� vznikne dek�dov�n�m genotypu). Evolu�n� strategie byly prvn�m p��stupem specializuj�c�m se na optimalizaci re�ln�ch parametr� a tak� p�inesly prvn� koncept meta-evoluce, optimalizace parametr� evoluce pomoc� vlastn�ho evolu�n�ho algoritmu. Genetick� programov�n� je p��kladem �sp�n� evoluce slo�it�ch struktur syntaktick�ch strom� po��ta�ov�ch program�. Neuroevoluce ukazuje mo�nosti vyu�it� evolu�n�ch technik pro ur�en� struktury i parametr� model� neuronov�ch s�t�. Oblast rojov�ch algoritm� se inspiruje chov�n�m spole�ensk�ho hmyzu a hejn pro efektivn� prohled�v�n� r�zn�ch prostor�, nap�. re�ln�ch euklidovsk�ch prostor� nebo hled�n� cest v grafech. 


\chapter{Evoluční algoritmy}
%!TEX root = main.tex
\index{Evolu�n� algoritmus}

%%%%%%%%%%%%%%%%%%%%%%%%%%%%%%%%%%%%%%%%%%%%%%%%%%%%%%%%%%
\section{Evoluce, geny a DNA}
%%%%%%%%%%%%%%%%%%%%%%%%%%%%%%%%%%%%%%%%%%%%%%%%%%%%%%%%%%


%%%%%%%%%%%%%%%%%%%%%%%%%%%%%%%%%%%%%%%%%%%%%%%%%%%%%%%%%%
\section{Obecn� sch�ma evolu�n�ho algoritmu}
%%%%%%%%%%%%%%%%%%%%%%%%%%%%%%%%%%%%%%%%%%%%%%%%%%%%%%%%%%

Oblast evolu�n�ch v�po�t� �i algoritm� (v angli�tin� evolutionary computing) zast�e�uje n�kolik proud�, kter� se zpo��tku vyv�jely samostatn�. Za prehistorii t�to discipl�ny lze pova�ovat Turingovy n�vrhy na vyu�it� evolu�n�ho prohled�v�n� z roku 1948~\cite{Turing:1948}\missingref a Bremermannovy prvn� pokusy o implementaci optimalizace pomoc� evoluce a rekombinace z roku 1962~\cite{Bremermann:1962}\missingref. B�hem �edes�t�ch let se objevily t�i skupiny v�zkumn�k�, kte�� nez�visle na sob� vyv�jely a navrhly sv� varianty pou�it� evolu�n�ch princip� v informatice. Holland publikoval v roce 1975 sv�j n�vrh genetick�ch algoritm�~\cite{Holland:1992}, \index{Genetick� algoritmus}
zat�mco skupina Fogela a spolupracovn�k� vyvinula metodu nazvanou evolu�n� programov�n�~\cite{Fogel:1995}. 
\index{Evolu�n� programov�n�}
Nez�visle na nich p�i�li Rechenberg a Schwefel v N�mecku na metodu nazvanou evolu�n� strategie~\cite{Beyer:2002}.
\index{Evolu�n� strategie}
A� do p�elomu osmdes�t�ch a devades�t�ch let existovaly tyto sm�ry bez v�razn�j�� interakce, ale pot� se spojily do obecn�j�� oblasti evolu�n�ch algoritm�. V t� dob�  Koza vytv��� metodu genetick�ho programov�n�, Dorigo publikuje disertaci s n�vrhem mraven��ch optimaliza�n�ch algoritm� a vznikaj� prvn� pokusy o aplikaci evoluce na v�voj um�l�ch neuronov�ch s�t�. 

U zrodu r�zn�ch variant evolu�n�ch algoritm� st�la inspirace p��rodn�mi fenom�ny, konkr�tn� jde o Darwinovu teorii p��rodn�ho v�b�ru a zjednodu�en� principy genetiky, kter� poprv� na�rtl Mendel. Z genetiky se evolu�n� algoritmy inspiruj� diskr�tn� reprezentac� genotypu, z biologick� evolu�n� teorie pou��vaj� Darwinovu my�lenku o v�b�ru jedinc� v prost�ed� s omezen�mi zdroji, kter� z�vis� na m��e p�izp�soben� se jedinc� dan�mu prost�ed�.

Z�kladn� obecnou my�lenku evolu�n�ch algoritm� lze vyj�d�it n�sleduj�c�m zp�sobem. M�jme populaci jedinc� v prost�ed�, kter� ur�uje jejich �sp�nost --- fitness. Tito jedinci navz�jem soupe�� o mo�nost reprodukce a p�e�it�, kter� z�vis� pr�v� na hodnot� fitness. Jde tedy o mno�inu kandid�t� na �e�eni probl�mu definovan�ho prost�ed�m. Zp�soby reprezentace jedinc�, jejich v�b�ru a rekombinace z�vis� na konkr�tn�m dialektu evolu�n�ch algoritm�, kter� probereme vz�p�t�. 

\begin{algorithm}
\caption{Sch�ma evolu�n�ho algoritmu}
\label{obreva}
\begin{algorithmic}
\Procedure{Evolu�n� algoritmus}{} 
\State $t \gets 0$
\State \emph{Inicializuj} populaci $P_t$ n�hodn� vygenerovan�mi jedinci
\State \emph{Ohodno�} jedince v populaci $P_t$
\While{neplat� \emph{krit�rium ukon�en�}}
\State 	vyber z $P_t$ rodi�e \emph{Rodi�ovskou selekc�}
\State 	\emph{Rekombinac�} rodi�� vzniknou potomci
\State 	\emph{Mutuj} potomky
\State 	\emph{Ohodno�} potomky
\State 	\emph{Enviroment�ln� selekc�} vyber $P_{t+1}$ z $P_t$ a potomk�
\State $t \gets t+1$
\EndWhile
%\Until{krit�rium ukon�en�}
\EndProcedure
\end{algorithmic}
\end{algorithm}


Z�kladn� princip fungov�n� evolu�n�ch algoritm� je tedy n�sleduj�c�~\cite{Eiben:2003}. Na za��tku algoritmu vygenerujeme (nej�ast�ji n�hodn�) prvn� inici�ln� populaci jedinc�. V�echny jedince v populaci ohonot�me ohodnocovac� funkc�. Hodnota t�to funkce ur�uje �anci v�b�ru jedinc� b�hem rodi�ovsk� selekce. Vybran� jedinci jsou potom rekombinov�ni pomoc� rekombina�n�ho oper�toru, kter� typicky ze dvou jedinc� vytv��� jednoho �i dva potomky, a pomoci oper�toru mutace, kter� typicky prov�d� drobn� zm�ny jednoho jedince. T�mto postupem si vytvo��me mno�inu nov�ch kandid�t� �e�en�, a tito nov� jedinci potom sout�� s p�vodn�mi jedinci o m�sto v nov� populaci. V�b�r jedinc� do nov� populace (tedy jak�si slit� rodi�� a potomk�) m� na starosti envioment�ln� selekce berouc� v �vahu fitness jedinc� a p�ipadn� dal�� ukazatele jako je nap��klad st��� jedinc�. T�m je vytvo�ena nov� generace a tento cyklus pokra�uje do spln�n� ur�it�ho krit�ria ukon�en�, co� je nej�ast�ji dostate�n� dobr� nejlep�� jedinec nebo p�edem ur�en� po�et generac�.



\chapter{Genetické algoritmy}
%%%%%%%%%%%%%%%%%%%%%%%%%%%%%%%%%%%%%%%%%%%%%%%%%%%%%%%%%%
\section{Genetick� algoritmy}
%%%%%%%%%%%%%%%%%%%%%%%%%%%%%%%%%%%%%%%%%%%%%%%%%%%%%%%%%%
\index{Genetick� algoritmus}

Genetick� algoritmy jsou asi nejzn�m�j�� sou��st� evolu�n�ch v�po�t� a v~r�zn�ch obm�n�ch se pou��vaj� hlavn� p�i �e�en� optimaliza�n�ch �loh. Je zaj�mav�, �e p�vodn� Hollandovou motivac� p�i n�vrhu genetick�ho algoritmu bylo studovat vlastnosti p��rodou inspirovan� adaptace~\cite{Holland:1992}. Velk� ��st p�vodn� literatury byla v�nov�na popisu princip�, jak genetick� algoritmus pracuje p�i hled�n� �e�en� �lohy. Zaj�mav� jsou paralely s matematick�m probl�mem dvouruk�ho bandity, kter� je p��kladem na udr�ov�n� optim�ln� rovnov�hy mezi explorac� a exploatac�. 

Nejjednodu��� varianta genetick�ho algoritmu pracuje s bin�rn�mi jedinci, to znamen�, �e parametry �e�en� �lohy je nutno v�dy zak�dovat jako bin�rn� �et�zce. Tento p��stup je v�hodn� z hlediska jednoduchosti pou�it�ch oper�tor�, ale bin�rn� zak�dov�n� na druhou stranu nemus� b�t nejvhodn�j�� reprezentac� probl�mu. Zp�sob fungov�n� jednoduch�ho genetick�ho algoritmu je tak� pom�rn� jednoduch�. Algoritmus p�ech�z� mezi populacemi �e�en� tak, �e nov� populace zcela nahrad� p�edchoz�. V�b�r rodi�� je �asto realizov�n tzv. ruletovou selekc�, kter� vyb�r� jedince n�hodn�~s pravd�podobnost� v�b�ru �m�rn� jejich fitness. Rekombina�n�m oper�torem je jednobodov� k���en�, kter� n�hodn� zvol� stejnou pozici v rodi��ch a vym��� jejich ��sti. Pravd�podobnost uskute�n�n� operace k���en� je jedn�m z parametr� programu a obvykle je pom�rn� vysok� (0,5 i v�ce). Mutace prov�d� drobn� lok�ln� zm�ny tak, �e proch�z� jednotliv� bity �et�zce a ka�d� bit s~velmi malou pravd�podobnost� zm�n�. Pravd�podobnost mutace je typicky nastavena, tak aby do�lo pr�m�rn� ke zm�n� jednoho bitu v populaci (obl�ben� doln� mez) nebo v jedinci (horn� mez). 

\begin{algorithm}
\caption{Sch�ma Hollandova gentick�ho algoritmu}
\label{obrga}
\begin{algorithmic}
\Procedure{Jednoduch� genetick� algoritmus}{} 
\State $t \gets 0$
\State \emph{Inicializuj} populaci $P_t$ $N$ n�hodn� vygenerovan�mi bin�rn�mi jedinci d�lky $n$
\State \emph{Ohodno�} jedince v populaci $P_t$
\While{neplat� \emph{krit�rium ukon�en�}}
\For{$i \gets 1, \dots, N/2$}
\State 	vyber z $P_t$ 2 rodi�e \emph{Ruletovou selekc�}
\State 	S pravd�podobnost� $p_C$ \emph{Zk�i�} rodi�e
\State 	S pravd�podobnost� $p_M$ \emph{Mutuj} potomky
\State 	\emph{Ohodno�} potomky
\State  P�idej potomky do $P_{t+1}$
\EndFor
\State 	Zaho� $P_t$
\State $t \gets t+1$
\EndWhile
\EndProcedure
\end{algorithmic}
\end{algorithm}

Ruletovou selekci si dle metafory m��eme p�edstavit tak, �e kolo rulety rozd�l�me na v�se�e odpov�daj�c� velikost� hodnot�m fitness jedinc� a p�i v�b�ru pak $n$ kr�t vhod�me kuli�ku. �asto pou��van�m vylep�en�m ruletov� selekce je tzv. stochastick� univerz�ln� v�b�r, kter� hod� kuli�ku do rulety jen jednou a dal�� jedince vyb�r� deterministicky posunem pozice kuli�ky o $1/n$. Tento v�b�r pro mal� $n$ l�pe aproximuje ide�ln� po�ty zastoupen� jedinc� v dal�� generaci. Dal��mi varianty rodi�ovsk� selekce nepracuj� s absolutn�mi hodnotami fitness, ale vyb�raj� n�hodn� v z�vislosti na po�ad� jedince v populaci set��d�n� podle fitness, co� zanedb�v� absolutn� rozd�ly mezi hodnotami. Dal�� variantou rodi�ovsk� selekce je tzv $k$-turnaj, kdy nejprve vybereme $k$ jedinc� n�hodn� a z nich pak vybereme nejlep��ho. 

\subsection{Oper�tory}

\subsection{GA na ��slech}

\subsection{GA na permutac�ch}



V sou�asnosti se oblast genetick�ch algoritm� neomezuje jen na bin�rn� k�dov�n� jedinc�, �ast� je celo��seln�, permuta�n� nebo re�ln� k�dov�n�, kter� ale vy�aduj� specifick� operace k���en� a mutace~\cite{Michalewicz:1996, Mitchell:1996}. O n�kter�ch se zm�n�me d�le. 


\chapter{Evoluční programování}
Evoluční programování je oblast, která je na rozdíl od genetických algoritmů charakterizována velmi benevolentním přístupem ke kódování jedinců. První varianty tohoto přístupu vznikly jako metoda pro automatickou evoluci konečných automatů \cite{fogel}. Jedincem v populaci je v tomto připadě tedy konečný automat (zakódovaný jakkoliv) a variační operace se omezují na několik různých mutací. Pro evoluční programování je charakteristické, že mutace jsou přizpůsobené problému a je jich často více, zatímco křížení se vůbec nepoužívá. 

\subsection{EP na KA}

\subsection{EP na číslech}

V dalším textu popíšeme variantu evolučního programování, která pracuje s jedinci reprezentovanými vektory reálných čísel. Rodičovská selekce je jednoduchá, každý rodič je jednou vybrán a mutací dá vznik jednomu potomkovi. Enviromentální selekce potom z množiny rodičů a potomků vybere turnajovou selekcí polovinu, která se stane další generací. 

\begin{algorithm}
\caption{Schéma meta-evolučního programování nad vektorem reálných čisel}
\label{obrep}
\begin{algorithmic}
\Procedure{Meta-EP}{} 
\State $t \gets 0$
\State \emph{Inicializuj} populaci $P_t$ $N$ náhodně vygenerovanými reálnými vektory $\vec{x^t}=(x_1^t,\dots,x_n^t,\sigma_1^t,\dots,\sigma_n^t)$
\State \emph{Ohodnoť} jedince v populaci $P_t$
\While{neplatí \emph{kritérium ukončení}}
\For{$i \gets 1, \dots, N$}
\State 	k rodiči $\vec{x_i^t}$ vygenruj potomka $\vec{y_i^t}$ \emph{mutací}:
\For{$j \gets 1, \dots, n$}
\State $\sigma'_j \gets \sigma_j \cdot (1+\alpha \cdot N(0,1))$
\State $x'_j \gets x_j + \sigma'_j \cdot N(0,1)$

\EndFor
\State vlož $\vec{y_i^t}$ do kandidátské populace potomků $P'_t$
\EndFor
\State 	\emph{Turnajovou selekcí} vyber $P_{t+1}$ z rodičů $P_t$ a potomků $P'_t$
\State 	Zahoď $P_t$ a $P'_t$
\State $t \gets t+1$
\EndWhile
\EndProcedure
\end{algorithmic}
\end{algorithm}
 
Mutace v našem připadě pracuje tak, že se snaží o malou změnu reálné hodnoty, kterou realizuje výběrem z normálního rozdělení. Ukázalo se, že hodnoty rozptylu normálního rozdělení jednotlivých parametrů jsou podstatné pro dobré fungování algoritmu. Jedno z navržených řešení je začlenit tyto rozptyly do každého jedince a upravovat je také v průběhu algoritmu. Jelikož tak vlastně upravujeme parametry, které zároveň používáme při vlastní mutaci, hovoříme o meta-evoluci.

\subsection{Význam křížení pro EA}


\chapter{Spojitá optimalizace}
%!TEX root = main.tex

Mnoho optimalizačních problémů z běžného života se dá definovat jako optimalizace funkce $$f: D \to \Real\,,$$ kde $D \subseteq \Real^d$. Je proto přirozené, že mnoho výzkumníků se zabývá právě evolučními algoritmy, kterou jsou schopné optimalizovat takové funkce. Problém optimalizace takových funkcí se v literatuře objevuje pod pojmem \emph{spojitá optimalizace}\footnote{anglicky \emph{continuous optimization}}. Je důležité si uvědomit, že název
 vyjadřuje pouze to, že prostor, ve kterém se hledají řešení je spojitý ($\Real^n$), samotná optimalizovaná funkce $f$ být spojitá nemusí.

Velmi často je definiční obor funkce $D$ $d$-rozměrný interval $$D = [l_1, u_1] \times \dots \times [l_d, u_d]\,,$$ kde $l_i$ a $u_i$ jsou dolní a horní meze pro $i$-tou proměnnou. Objevují se ale i obecnější problémy, kde je $D$ dána pomocí množiny podmínek tvaru $g(\vec{x}) \leq 0$ a $h(\vec{x}) = 0$, pro $g, h: \Real^d \to \Real$.

V této kapitole napřed budeme uvažovat optimalizační problém tvaru 
\begin{equation*}
  \begin{aligned}
  & \underset{x}{\text{min}} & & f(\vec{x}) \\
  & \text{za podmínek} & & \vec{x} \in [l_1, u_1] \times \dots \times [l_d, u_d]\,. \\
  \end{aligned}
\end{equation*}

\section{Vlastnosti funkcí}

Je zřejmé, že některé typy funkcí budou pro evoluční algoritmy lehčí, než jiné. Velký vliv na efektivitu evolučního algoritmu mají především vlastnosti jako multi-modalita, separabilita a podmíněnost. 

Funkce je \emph{multi-modální}, pokud má velké množství lokálních optim. Je zřejmé, že v takovém případě může mít algoritmus problém s uváznutím v lokálním optimu a je potřeba tomu přizpůsobit operátory. Existuje i oblast multi-modální optimalizace, kde je cílem najít co nejvíce různorodých lokálních optim. 

Separabilní funkce jsou naopak pro optimalizaci jednodušší. Jsou to takové funkce, které se dají zapsat pomocí funkcí jedné proměnné. Formálně, funkce $f(x_1, \dots, x_n): \Real^n \to \Real$ je \emph{aditivně separabilní}, pokud se dá zapsat jako součet funkcí $f_1(x_1), \dots, f_n(x_n)$, tj. $f(x_1, \dots, x_n)= \sum_{i=1}^n f(x_i)\,.$ Obdobně můžeme zadefinovat i funkci multiplikativně separabilní. Z hlediska optimalizace je velkou výhodou separabilních funkcí, že se dají optimalizovat po jednotlivých složkách vektoru, tj. optimum můžeme najít tak, že vždy zafixujeme hodnoty $n-1$ parametrů a optimalizujeme jen podle jednoho.

Další vlastností, které výrazně ovlivňuje evoluci je podmíněnost funkce. Ta vyjadřuje, jak moc se liší vliv jednotlivých proměnných na hodnotu funkce. Pro kvadratické funkce (u kterých vrstevnice vypadají jako elipsoidy), je jejich podmíněnost druhá odmocnina poměru mezi délkou nejdelší a nejkratší osy elipsoidu. Pokud je podmíněnost funkce velká, říká se, že je funkce špatně podmíněná. Z hlediska evolučního algoritmu je důležité, že by s různým vlivem proměnných na hodnotu funkce měly počítat operátory.

\begin{marginfigure}
\centering
%!TEX root = ../../main.tex
\subfigure[Dob�e podm�n�n� separabiln� funkce]{
  \begin{tikzpicture}
  \begin{axis}[
    axis x line=center,
    axis y line=center,
    xmin=-5,   xmax=5,
    ymin=-5,   ymax=5,
    xlabel=$x_1$,
    ylabel=$x_2$,
    width=\marginparwidth
  ]
    \pgfplotsinvokeforeach{0.2, 0.4, 0.6, 0.8, 1.0}
    {\draw[red] \pgfextra{
      \pgfpathellipse{\pgfplotspointaxisxy{0}{0}}
      {\pgfplotspointaxisdirectionxy{#1*4.5}{0}}
      {\pgfplotspointaxisdirectionxy{0}{#1*4.5}}
    };}
  \end{axis}
  \end{tikzpicture}
}
\subfigure[�patn� podm�n�n� funkce]{
  \begin{tikzpicture}
  \begin{axis}[
    axis x line=center,
    axis y line=center,
    xmin=-5,   xmax=5,
    ymin=-5,   ymax=5,
    xlabel=$x_1$,
    ylabel=$x_2$,
    width=\marginparwidth
  ]
    \pgfplotsinvokeforeach{0.2, 0.4, 0.6, 0.8, 1.0}
    {\draw[red] \pgfextra{
      \pgfpathellipse{\pgfplotspointaxisxy{0}{0}}
      {\pgfplotspointaxisdirectionxy{#1*4.5}{0}}
      {\pgfplotspointaxisdirectionxy{0}{#1}}
    };}
  \end{axis}
  \end{tikzpicture}
}
\subfigure[�patn� podm�n�n� a neseparabiln� funkce]{
  \begin{tikzpicture}
  \begin{axis}[
    axis x line=center,
    axis y line=center,
    xmin=-5,   xmax=5,
    ymin=-5,   ymax=5,
    xlabel=$x_1$,
    ylabel=$x_2$,
    width=\marginparwidth
  ]
    \pgfplotsinvokeforeach{0.2, 0.4, 0.6, 0.8, 1.0}
    {\draw[red] \pgfextra{
      \pgfpathellipse{\pgfplotspointaxisxy{0}{0}}
      {\pgfplotspointaxisdirectionxy{#1*4.5}{#1*1.8}}
      {\pgfplotspointaxisdirectionxy{0}{#1}}
    };}
  \end{axis}
  \end{tikzpicture}
}
\caption{Příklady různých vlastností funkcí}
\label{fig:function_examples}
\end{marginfigure}

Příklady posledních dvou vlastností vidíme na obrázku \ref{fig:function_examples}, který ukazuje vrstevnice funkcí dvou proměnných. Horní funkce je jednoduchý dvoudimenzionální paraboloid, který je separabilní a dobře podmíněný. Na prostředním obrázku je paraboloid, který má jednu osu delší než druhou a znázorňuje tedy špatnou podmíněnost\footnote{podmíněnost zobrazené funkce jen cca 4.8, nedá se tedy považovat za špatně podmíněnou, objevují se i funkce s podmíněností $10^6$}. Poslední funkce potom kombinuje špatnou podmíněnost s neseparabilitou.

Výše uvedené vlastnosti jsou ty, které nejvíce ovlivňují efektivitu evolučních algoritmů při spojité optimalizaci, nejsou to ale všechny. Některé algoritmy například mohou využívat různé symetrie dané funkce, naopak funkce, které mají své globální optimum jen ve velmi malé oblasti prohledávaného prostoru a jinak jsou konstantní jsou pro evoluci velmi těžké obecně.

\section{Kódování pro spojitou optimalizaci}

Jedno z prvních rozhodnutí, které je potřeba udělat při návrhu evolučního algoritmu je výběr kódování jedince. Vzhledem k tomu, že ve spojité optimalizaci pracujeme s vektory reálných čísel, je otázka výběru kódování relativně snadná a jedinci jsou v naprosté většině případů kódování jako vektor typu \code{float} nebo \code{double}.


\ask{Existuje pro tohle $\downarrow$ nějaká reference?}

Dalo by se uvažovat i o kódování jedince přímo po vzoru Hollandova genetického algoritmu, tj. binárním vektorem, a používat jednoduchá $n$-bodová křížení a bit-flip mutace, ale taková reprezentace trpí tím, že změna různých bitů v číslech vede k výrazně různým změnám hodnoty (např. změna bitu na konci mantisy vs. změna bitu na začátku exponentu).

\section{Operátory pro spojitou optimalizaci}

Pro spojitou optimalizaci můžeme samozřejmě použít stejné operátory jako pro každé jiné vektorové kódování jedince, tedy například $n$-bodové, případně uniformní křížení. Nicméně častěji se používají operátory specializované, které přímo využívají toho, že jedinci jsou vektory čísel.

Tzv. aritmetické křížení vektorů počítá vážený průměr dvou rodičů tak, aby vytvořilo potomka, potomci se tedy spočítají jako
\begin{align*}
o_1 & = w*p_1 + (1-w)* p_2\,, \\
o_2 & = (1-w)*p_1 + w* p_2\,,
\end{align*}
kde $p_1$ a $p_2$ jsou rodiče a $w \in (0,1)$ je (případně náhodně) zvolená váha. Ačkoliv se takové křížení relativně často používá hlavně v jednodušších aplikacích, jeho velkou nevýhodou je, že výslední potomci se nikdy nemohou dostat z konvexního obalu počáteční populace.

Mutace pro spojitou optimalizaci se občas rozdělují na dva typy --- ovlivněné a neovlivněné. Neovlivněná mutace generuje novou hodnotu pro složku vektoru nezávisle na aktuální hodnotě, ovlivněná mutace naopak aktuální hodnotu používá. Typickým příkladem neovlivněné mutace je vygenerování nového čísla z daného rozsahu. Ovlivněná mutace typicky přičítá k dané složce vektoru číslo z normálního rozdělení $\normal{\mu}{\sigma^2}$.

Jednou z nevýhod výše uvedeného aritmetického křížení je, že potomci jsou velmi často relativně daleko od svých rodičů a nejsou jim tedy moc podobní. Právě z toho důvodu Deb a Agrawal navrhli simulované binární křížení (SBX)\cite{agrawal1995simulated}.

\chapter{Evoluční strategie}
%!TEX root = main.tex

\emph{Evoluční strategie}\cite{rechenberg1973,schwefel1977numerische} jsou z historického hlediska jistou alternativou Hollandova genetického algoritmu. Jsou o něco starší a je na nich zajímavé, že jsou o něco komplikovanější. V současnosti se používají především pro řešení problémů spojité optimalizace, ale myšlenky, které se v této oblasti vyskytují, se dají použít i jinde.

Evoluční strategie se dělí do dvou skupin podle způsobu, jakým pracují s populacemi rodičů a potomků. V obou případech jsou důležité parametry $\mu$ a $\lambda$, které označují počet rodičů a počet potomků, kteří z nich vznikají. Pro druhy evolučních strategií potom existuje ustálené značení $(\mu, \lambda)$-ES a $(\mu + \lambda)$-ES. V prvním případě (``čárková selekce'') máme populaci $\mu$ rodičů, ze kterých vytvoříme $\lambda$ potomků ($\lambda > \mu$), z těch potom vybereme nejlepších $\mu$ jako rodiče do další generace.  Ve druhém případě (``plus selekce'') z $\mu$ rodičů vytvoříme opět $\lambda$ potomků, ale před selekcí napřed sloučíme rodiče a potomky do jedné populace velikosti $\mu + \lambda$. Z té se potom opět vybere $\mu$ nejlepších jedinců jako rodiče do další generace. 

Jednou ze základních vlastností evolučních strategií, které je odlišují od jiných typů evolučních algoritmů je to, že obsahují nějakou formu samo-adaptace parametrů. V případě spojité optimalizace se tedy kromě samotných hodnot vektoru vyvíjí například i parametry pro mutaci. Technicky se tedy potom jedinec skládá ze dvou částí -- samotného zakódovaného vektoru čísel $\vec{x}$ a vektoru tzv. \emph{endogenních parametrů} $\vec{s}$, které právě obsahují všechny parametry, které ovlivňují chování operátorů\footnote{Vedle pojmu ``endogenní parametry'' se ještě občas objevuje pojem ``exogenní parametry'', který označuje parametry algoritmu, které se nemění, jako např. velikost populace.}. Je důležité si uvědomit, že endogenní parametry nijak přímo neovlivňují fitness jedince, jejich hodnoty se vyvíjí jen díky tomu, že jedinci s lepší hodnotou endogenních parametrů mají po aplikaci genetických operátorů častěji lepší fitness. Pro samotnou evoluci endogenních parametrů se mohou používat stejné operátory jako pro samotného jedince s fixně nastavenými parametry, nicméně moderní evoluční strategie častěji používají deterministický způsob nastavení těchto parametrů.

Důležitou součástí evolučních strategií je rekombinace. Ta v zásadě odpovídá křížení v genetických algoritmech a jejím cílem je vytvořit nového jedince kombinací jedinců z populace. V evolučních strategiích se ale velmi často používají rekombinace, které kombinují všechny jedince. Častá je tzv. \emph{interpolační rekombinace}, která jednoduše spočítá průměr všech jedinců v populaci -- nový jedinec je vlastně centroid celé populace. Používá se i její varianta, která používá vážený součet, kde lepší jedinci v populaci mají větší váhu při výpočtu průměru (nejhorší jedinci se dokonce mohou úplně ignorovat). Kromě těchto dvou nejčastěji používaných rekombinací se dají používat i křížení známá z genetických algoritmů, není to ale moc obvyklé. Existuje i uniformní rekombinace, která každou složku vektoru jedince vybírá z náhodného jedince v populace. 

Použitá rekombinace se občas objevuje i v notaci evolučních strategií, v takovém případě se píše jako $(\lambda/\rho + \mu)$, kde právě $\rho$ označuje použitou rekombinaci.

Základním operátorem v evolučních strategiích ale je mutace. Typicky jde o gaussovskou mutaci s určitým rozptylem. Zde se objevuje několik možností, kde nejjednodušší je sférická mutace, která přičítá k jedinci vektor z normálního rozdělení $\sigma\normal{0}{I}$, kde $I$ je jednotková matice. O něco složitější varianta potom používá jiný rozptyl pro každou souřadnici vektoru, tj. k jedinci se přičítá hodnota z normálního rozdělení $\normal{0}{\mathrm{diag}({\vec{\sigma}})}$, kde $\mathrm{diag}(\sigma)$ značí diagonální matici s vektorem $\vec{\sigma}$ na hlavní diagonále. Konečně v nejkomplikovanějším případě se jedinci generují s normálního rozdělení s plnou kovarianční matici $\normal{0}{\Sigma}$. Ačkoliv první dva způsoby vyžadují mnohem nižší množství endogenních parametrů (1 respektive $d$), jejich nevýhoda spočívá v tom, že nejsou schopny zachytit závislosti mezi parametry a hodí se tak především pro separabilní problémy. Na druhou stranu způsob s plnou kovarianční maticí je invariantní vůči rotacím a škálovaní prohledávaného prostoru (a díky dalším vlastnostem evolučních strategií i k monotónnímu škálování fitness funkce).

Vzhledem k relativně velkému množství endogenních parametrů v evolučních strategiích a k tomu, že vhodné nastavení parametrů je různé pro různé optimalizační problémy a dokonce i v různých fázích evoluce, vznikla potřeba parametry nastavovat adaptivně. Při pevném nastavení velikosti mutace typicky pozorujeme tři fáze evoluce. V první fázi evoluce je konvergence pomalá, protože změny dělané mutací jsou příliš malé. Následuje fáze rychlé konvergence, kde jsou velikosti změn optimální pro danou fázi výpočtu, a v poslední fázi se konvergence zase zpomalí, protože jsou změny moc velké (algoritmus ``skáče'' kolem optima).

\newcommand{\fifth}{\sfrac{1}{5}\xspace}

\section{Pravidlo jedné pětiny}

První metoda adaptivního nastavování vzešla z Rechenbergových experimentů,\cite{rechenberg1973} které zkoumaly vliv rozptylu Gaussovské mutace v $(1+1)$-ES při optimalizaci dvou funkcí. Ukázalo se, že pro obě funkce algoritmus nejrychleji konverguje, pokud se velikost mutace zmenší, když je pravděpodobnost, že je potomek lepší než rodič menší než cca \fifth, a zvětší, když je tato pravděpodobnost větší než cca \fifth. Z tohoto pozorování vzniklo pravidlo \fifth.\sidenote{V angličtině ``one fifth rule''}

Proč takové pravidlo funguje? Představme si optimalizaci lineární funkce, v takovém případě je pravděpodobnost, že potomek je lepší než rodič, přesně 50\%. 
Obecnou funkci můžeme v každém bodě podle Taylorova pravidla aproximovat pomocí lineární funkce, tato aproximace je tím přesnější, čím blíže jsme bodu, ve kterém jsme funkci aproximovali, tedy, při malé velikosti mutace bude většina potomků blízko a pravděpodobnost, že jsou lepší než rodič, je kolem 50\%. Naopak, při nekonečném rozptylu mutace pravděpodobnost zlepšení odpovídá části prostoru, kde má funkce lepší hodnotu, než v daném bodě. Ve chvíli, kdy se algoritmu přiblíží optimu je tato pravděpodobnost většinou blízko 0. Ve skutečnosti se dokonce ukazuje, že pro většinu funkcí pravděpodobnost zlepšení monotónně roste s klesajícím rozptylem mutace. Přesná hodnota pro optimální pravděpodobnost zlepšení závisí na optimalizované funkci, nicméně právě \fifth je často doporučovaná.

\section{Sebe-adaptivní evoluční strategie}

Jiný přístup k sebe-adaptaci endogenních parametrů zvolil Schweffel\cite{schwefel1977numerische}, ten pro jejich ladění použil podobné genetické operátory, jako se používají pro samotné zakódované řešení. Konkrétně se každý endogenní parametr vynásobil číslem $e^{\normal{0}{I}}$.

\chapter{Kombinatorická optimalizace}
%\section{Kombinatorická optimalizace}
%
%

\section{Problém obchodního cestujícího}

Problém obchodního cestujícího (POC) je jedním z nejznámějších NP-úplných problémů. Jde o nalezení nejkratší cesty mezi $n$ městy, která navštíví všechna města právě jednou a skončí v počátečním městě~\footnote{V daném ohodnoceném úplném grafu máme najít nejkratší hamiltonovskou kružnici.}. 

Evoluční řešení problému obchodního cestujícího má zřejmou fitness --- délku cesty --- zatímco reprezentací řešení existuje několik a genetické operátory (hlavně operátor křížení) jsou na těchto reprezentacích velmi závislé. V následujícím ukážeme nejprve nejčastější reprezentaci pomocí cesty (permutace)~\footnote{Permutační zakódování jedinců je přirozené i u jiných problémů, zapamatujme si jej.}, které je sice velmi přímočará, ale vyžaduje specifické operátory křížení. O operátorech mutace a inicializace se zmíníme poté.   


\subsection{Reprezentace cestou}

Reprezentace cestou je velmi názorná, jde o posloupnost celočíselných hodnot označujících pořadí měst tak, jak je navštívíme. Jedinci evolučního algoritmu tedy permutaci o $n$ prvcích.  $(412876935)$ tedy říká, že řešení reprezentuje cestu městy v pořadí $4-1-2-8-7-6-9-3-5$. Zjevnou nevýhodou této reprezentace je, že jednobodové křížení dvou permutací typicky nevyrobí potomky jako validní permutace. 

Jednotlivé návrhy křížení proto pracují tak, aby opravily nebo doplnily jedince, který od svých rodičů zdědí vlastnosti, jež autoři považují za důležité. V různých kříženích se proto (kromě zachování permutace) zrcadlí heuristický přístup k tomu, co by měla rekombinace dvou cest v problému POC zohlednit. 

Prvním návrhem křížení je PMX --- křížení částečné shody nebo také částečného zobrazení.\footnote{Partially mapped/matched crossover} Goldberga a Linglea. Snahou tohoto křížení je zachovat co nejvíce měst na svých původních pozicích. Jde o dvoubodové křížení, které dvěma náhodnými řezy rozdělí každého rodiče na tři části, přičemž střední část použije k definici zobrazení, která města na stejných pozicích spolu korespondují. V příkladu na obr.~\ref{PMX} definujeme tedy zobrazení měst: $1\leftrightarrow 4, 8\leftrightarrow 5, 7\leftrightarrow 6, 6\leftrightarrow 7$. Křížení PMX vymění středy jedinců a dále doplní z původních cest ta města, která nezpůsobí konflikt, t.j. nejsou ve vyměněných středních částech. Nakonec se pro doplnění zbylých měst použije korespondence dané zobrazením dle střední části.   

\begin{figure}%
\begin{verbatim}
(123|4567|89) PMX (452|1876|93) :
(...|1876|..)  (...|4567|..) 
a zobrazení: 1-4 8-5 7-6 6-7 
lze doplnit (.23|1876|.9)  (..2|4567|93)
a dle zobr. (423|1876|59)  (182|4567|93)
\end{verbatim}
\caption{Příklad křížení PMX.}%
\label{PMX}%
\end{figure}

Davis navrhl jiný přístup ke křížení, tzv. křížení zachovávající pořadí (OX)~\footnote{Order crossover}, které se snaží co nejvíce zachovat relativní pořadí jednotlivých měst v cestě. Obdobně jako v případě PMX jde o dvoubodové křížení, které nejprve vymění středy jedinců. Dále přeuspořádáme cestu druhého rodiče od druhého bodu křížení a popořadě doplňujeme prvního potomka tak, že vynecháváme města, která již jsou v nové střední části. Druhý potomek vznikne stejným způsobem z přeuspořádané cesty prvního rodiče, viz obr.~\ref{OX}.   


\begin{figure}%
\begin{verbatim}
(123|4567|89) OX (452|1876|93) : 
(...|1876|..)  (...|4567|..) 
a přeuspořádáme cestu 2 od druhého bodu 
křížení: 9-3-4-5-2-1-8-7-6
vyhodíme překřížená města z 1, zbyde: 9-3-2-1-8
doplníme potomka 1: (218|4567|93)
obdobně potomek 2: (345|1876|92)
\end{verbatim}
\caption{Příklad křížení OX.}%
\label{OX}%
\end{figure}

Na základě křížení OX vznikl algoritmus křížení modifikovaného křížení (ModX)~\footnote{Modified crossover}, které se ukázalo jako úspěšné v řešení problémů rozvrhování. ModX je vlastně OX jen s jedním křížícím bodem, kde oba rodiče přeuspořádáme tak, aby cesta začínala za křížícím bodem a pak křížíme dle principu OX od začátku obou jedinců, viz obr.~\ref{ModX}.    

\begin{figure}%
\begin{verbatim}
(0845|671239) ModX (6712|483590)

potomek 1: (6712084539)
potomek 2: (0845671239) 
\end{verbatim}
\caption{Příklad křížení ModX.}%
\label{ModX}%
\end{figure}

Posledním z klasických křížení permutačních řešení POC, které si uvedeme, je cyklické křížení (CX)~\footnote{Cycle crossover} navržené Davisem a Hollandem. Snahou autorů bylo, aby každá pozice v potomkovi odpovídala, pokud možno, jednomu s rodičů. Za tímto účelem se při tvorbě potomka pohybujeme v cyklech po permutaci podle odpovídajících měst. První pozici si zvolíme náhodně z jednoho z rodičů, pak se snažíme vzít to město, které je na stejné pozici u druhého z rodičů, ale bereme z pozice v prvním rodiči, atd., dokud nedokončíme cyklus. Pak doplníme městy a pozicemi z druhého rodiče, viz obr.~\ref{CX}.  

\begin{figure}%
\begin{verbatim}
(123456789) CX (412876935)
první pozice, náhodně třeba z prvního P1=(1........), 
teď musíme vzít 4, P1=(1..4....), pak 8, 3 a 2
P1=(1234...8.), dál už nelze, doplníme z druhého
P1=(123476985)
Obdobně P2=(412856739)
\end{verbatim}
\caption{Příklad křížení CX.}%
\label{CX}%
\end{figure}


\subsection{Křížení rekombinací hran}

Hlavní myšlenkou křížení rekombinací hran (ER)~\footnote{Edge recombination crossover} je přesun pozornosti od uzlů grafu (městům) k hranám (cestám). Autoři Whitley s kolegy zjistili, že výše zmíněné operátory ve skutečnosti nezachovávají více jak 60\% hran z obou rodičů. Algoritmus ER (a jeho vylepšení ER2) se tedy snaží pracovat tak, že v potomcích zachová co nejvíce již existujících hran. Pracuje ve dvou fázích --- nejprve si vybuduje seznam sousedů všech vrcholů, a potom tvoří nové jedince s pomocí tohoto seznamu, viz alg.~\ref{alger}.

\begin{algorithm}
\caption{Křížení rekombinací hran}
\label{alger}
\begin{algorithmic}
\Procedure{ER}{} 
\State $K \gets []$ prázdný seznam
\State $N \gets $ první uzel náhodného rodiče
\While{Length(K) $<$ Length(Parent)}
\State $K \gets [K,N]$ 
\State Vyhoď $N$ ze všech seznamů sousedů
\If{seznam sousedů $N \not\=\emptyset$} 
\State 	$N^* \gets$ soused s nejméně sousedy v seznamu (nebo náhodný z více takových)
\Else 
\State 	$N^* \gets$ náhodně zvolený uzel $\notin K$
\EndIf
\State 	$N \gets N^*$
\EndWhile
\EndProcedure
\end{algorithmic}
\end{algorithm}


Obr.~\ref{ER} a~\ref{ER2} ukazují práci původní varianty algoritmu a jejího vylepšení, které dává přednost výběru měst, která jsou v seznamu dvakrát (označené \#). Dodejme, že algoritmus se může zastavit ještě před sestavením kompletní permutace, ale to se dle praktických zkušeností děje jen zřídka (v 1--1,5\% případů). Algoritmy křížení ER jsou v praxi velmi úspěšné. 

\begin{figure}%
\begin{verbatim}
(123456789) ER (412876935)
1:   9 2 4
2:   1 3 8
3:   2 4 9 5
4:   3 5 1
5:   4 6 3
6:   5 7 9
7:   6 8
8:   7 9 2
9:   8 1 6 3

Začnu v 1, následníci jsou 9, 2, 4
9 vypadává, má 4 násl., z 2  a 4 náhodně vyberu 4
násl. 4 jsou 3 a 5, beru 5, 
neboť 5 má 3 násl, zatímco 3 má 4 násl. 
Teď máme (145......), podobně pokračujeme
... (145678239)
\end{verbatim}
\caption{Příklad křížení ER.}%
\label{ER}%
\end{figure}

\begin{figure}%
\begin{verbatim}
(123456789) ER (412876935)
1:   9 #2 4
2:   #1 3 8
3:   2 4 9 5
4:   3 #5 1
5:   #4 6 3
6:   5 #7 9
7:   #6 #8
8:   #7 9 2
9:   8 1 6 3

(12.......), (128......), (1287.....), (12876....)
náhodně např. (128765...), (1287654..), (12876543.)
(128765439)
\end{verbatim}
\caption{Příklad křížení ER2.}%
\label{ER2}%
\end{figure}

\subsection{Další operátory}

Zatím jsme hovořili jen o kříženích, zastavme se nyní u dalších operátorů pro problém POC. 

Inicializaci jedinců na začátku algoritmu můžeme provést generováním náhodných permutací, ale lze využít i heuristik pro lepší výchozí řešení. Jednou z nich je jednoduchý algoritmus hladového typu, který inicializuje tak, začne náhodným městem a pak postupně vybírá nejbližšího souseda k již určenému městu. Jinou možností inicializace je algoritmus vkládání hran~viz alg.~\ref{alginr}.


\begin{algorithm}
\caption{Inicializace vkládáním hran}
\label{alginr}
\begin{algorithmic}
\Procedure{Vkládání hran}{} 
\State cesta $T \gets$ náhodná hrana
\While{Length(K) $<n$}
\State vyber nejbližší město $c$ k cestě $T$
\State najdi hranu $(k-j)\in T$ tak, že minimalizuje rozdíl mezi $(k-c-j)$ a $(k-j)$
\State Vyhoď $(k-j)$ z $T$
\State Vlož $(k-c)$ do $T$
\State Vlož $(c-j)$ do $T$
\EndWhile
\EndProcedure
\end{algorithmic}
\end{algorithm}

Mutačních operátorů vhodných pro POC je celá řada, i když nejjednodušší (změna jednoho města) naopak použít nelze. Oblíbené jsou ale mutace 
Posun podcesty, Záměna 2 měst, Záměna podcest nebo heuristiky jako je 2-opt~viz alg.~\ref{dvaopt}~\footnote{G. A. CROES (1958). A method for solving traveling salesman problems. Operations Res. 6 (1958) , pp., 791-812.}. 
Zajímavé na permutační reprezentaci také je, že u ní funguje inverze. 


\begin{algorithm}
\caption{Heuristika 2-opt pro lokální zlepšení cesty}
\label{dvaopt}
\begin{algorithmic}
\Procedure{2-opt}{cesta, i, k} 
\State new1 $\gets$ cesta[1] .. cesta[i-1]
\State new2 $\gets$ cesta[i] .. cesta[k] v převráceném pořadí
\State new3 $\gets$ cesta[k+1] .. cesta[n]
\Return new1+new2+new3
\EndProcedure
\end{algorithmic}
\end{algorithm}


\subsection{Jiné reprezentace}

Přesto, že permutační zakódování jedinců POC je nejčastější, zkoumají se i jiné možnosti reprezentace cest. 

Jednou z nich je reprezentace sousednosti~\footnote{Adjacency representation}, která cestu reprezentuje také jako seznam měst, ale 
platí že město $j$ je na pozici $i$, vede-li cesta z $i$ do $j$. Takže např. 
(248397156) odpovídá cestě 1-2-4-3-8-5-9-6-7. Každá cesta má jen 1 reprezentaci, ale některé seznamy negenerují přípustnou cestu
Klasické křížení v této reprezentaci také nefunguje (a navrhují se různé opravy), ale hlavní motivací pro její zavedení bylo, že 
v ní mají dobrý význam klasická schémata. Např  (*3*...) značí všechny cesty s hranou 2-3

Problém fungování klasického jednobodového křížení se snaží řešit ordinální reprezentace~\footnote{reprezentace bufferem}. Ta je založena na myšlence, že číslo reprezentující jedno město nebude označovat jeho absolutní název ale jen relativní pozici (pořadí) v bufferu všech měst, ze kterého při konstrukci jedince použitá města mažeme. Takže například máme-li buffer (123456789) a cestu 1-2-4-3-8-5-9-6-7, její reprezentace je (112141311).
Relativita ordinální reprezentace způsobí, že nyní můžeme dva jedince křížit jednobodově (potomci jsou syntakticky správně), ale význam takového křížení není zřejmý, jde spíše o makromutaci.

Existuje i několik pokusů o  maticové reprezentace cest jako jedinců evolučního algoritmu, většinou jde o binární matice, kde jednička na pozici $(i,j)$ znamená buď, že v cestě vede hrana $i\to j$, anebo (častěji), že město $i$ předchází v cestě město $j$. Operátory nad maticemi používají často logické operace nad prvky matic a opravné mechanismy. Většinou tyto reprezentace nedosahují výrazných výsledků. 

\section{Jiné problémy}

\subsection{Batoh}


Evoluční algoritmy jsou aplikovány i na řešení dalších kombinatorických problémů, přičemž je zajímavé, že některé problémy, ač jsou např. také NP-těžké jako POC, představují jiné překážky pro EVA. Uvažme například problém 0-1 batohu, který je definován následovně: 

Mějme množinu věcí očíslovaných od 1 do $n$, každá má hmotnost $w_i$a cenu $v_i$, a maximální nosnost batohu W. Úkolem je
maximalizovat: $$\sum _{i=1}^{n}v_{i}x_{i}$$ za podmínek $$\sum _{i=1}^{n}w_{i}x_{i}\leq W\;\;\; \mbox{a zároveň}\;\;\; x_{i}\in \{0,1\}$$.
$x_i$ representuje počet věci $i$ v batohu, my počítáme s tím, že každou věc vezmeme maximálně jednou, proto 0-1 batoh. Zakódování problému se zdá --- na rozdíl od POC --- triviální, můžeme použít bitovou mapu pro věci určené do batohu. Operátory křížení a mutace se zdají být také jednoduché. Problémem v tomto případě je fitness jedince. Za prvé chceme maximalizovat cenu, za druhé nechceme překročit kapacitu $W$, ale blížit se k ní. Lze využít multi-kriteriální EVA, ale v praxi se osvědčilo i velmi jednoduché řešení z rodiny opravných algoritmů dekódování jedince, které bývá překvapivě efektivní. Dekodér procházející bitový vektor jedince interpretuje jedničku na místě $i$ jako: dej věc $i$ do batohu, pokud by se nepřeplní. 

\subsection{Rozvrhování}

Rozvrhování je dalším těžkým problémem, kde se evoluční algoritmy osvědčují jako heuristika pro nalezení přibližných řešení. Rozvrh reprezentujeme přímočaře maticí, kde řádky odpovídají učitelům, sloupce hodinám a hodnoty jsou kódy předmětů. Mutace většinou míchají předměty, křížení mění lepší řádky z jedinců. Problémem u rozvrhování je, jak určit fitness a jak pracovat s omezujícími podmínkami. Obecně se ukazuje, že dobré je definovat fitness řádku (t.j. jak je spokojen učitel) a tu pak zkombinovat s dalšími (často měkkými) kritérii kvality rozvrhu. Existující tvrdá omezení (kdy mohou učitelé, co kde lze učit, $\dots$) je z důvodů efektivity často nutné řešit již v operátorech, aby se zbytečně negenerovala řada nepřípustných řešení.   

\subsection{Plánování výroby}
Na problému plánování výroby~\footnote{shop scheduling} si můžeme ukázat, že různá zakódování nám vyřeší různé těžkosti s řešením (a přinesou jiná). Problém je definován takto: Mějme výrobky $o_1,\dots o_N$, sestávající z částí $p_1, \dots p_K$. Pro každou část máme stejný (flow shop) nebo více (job shop) plánů, jak jí vyrobit na strojích $m_1, \dots m_M$, a stroje mají různé doby nastavení na jiný výrobek. Úkolem samozřejmě je minimalizovat čas výroby. Tento čas se také jednoduše může stát fitness. 

Evoluční řešení volí dvě cesty jak zakódovat jedince, První typ kódování je jednoduchý permutační --– vyvíjený plán je jen permutace pořadí výrobků. Dekodér pak musí zvolit plány pro jejich části. Výhody jsou jednoduchá reprezentace a možnost použít křížení např. z problému POC. Ale ukazuje se, že tento přístup není efektivní. Dekodér je nucen řešit tu komplikovanou část úlohy, výběr plánu, a také většina operátorů z POC není vhodná.
Na druhou stranu, přímá reprezentace jedince jako celého plánu (nazývá se paralelní kódování strojů) klade nároky na specializované (komplexní) evoluční operátory~\cite{Werner:2013}.


\chapter{Evoluce pro strojové učení}
\section{Evoluce pravidlových systémů}

S rozvojem expertních systémů v 70. a 80. letech přišly i pokusy jak takové systémy znalostí navrhovat jinak než ručně ve spolupráci odborníků na umělou inteligenci s experty v dané problémové doméně. Znalosti v takových systémech, kterým se říká pravidlové, jsou většinou realizovány ve formě pravidel: 

IF (podmínka) THEN (výsledek)

Problém tedy je, jak navrhnout evoluční algoritmus, který bude vyvíjet daná pravidla nebo jejich množiny. Důležitým faktorem, který ovlivňuje návrh takového algoritmu, je typ učení, pro který ho chceme použít. Evoluce pravidlových systémů se historicky uvažuje v kontextu učení s učitelem (např. úlohy klasifikace) a zpětnovazebného učení (vývoj řídících mechanismů pro agenty či roboty). 

Obecně se dá říci, že pro vývoj pravidel v úlohách učení s učitelem máme jednodušší způsob výpočtu a práce s fitness, zatímco v úlohách zpětnovazebného učení (někdy též nazývané on-line učení) je mechanismus interakce pravidlového systému s prostředím a distribuce odměn mezi pravidly složitější.  

Základní principy pro evoluci pravidlových systémů byly položeny v 80.letech 20.století v rámci tzv. Michiganského a Pittsburghského přístupu. Od poloviny 90.let pak mluvíme o nové generaci klasifikačních systémů reprezentované algoritmem Rozšířený klasifikační systém (XCS). 


\section{Michiganský a Pittsburghský přístup}

Prvním, kdo se pokusil skloubit pravidlové systémy a evoluční algoritmy, byl J. Holland, který navrhl koncept \emph{učících se pravidlových systémů}. Pro Hollandův přístup je typické, že uvažoval od počátku problémy zpětnovazebného učení a jeho důraz je tedy na mechanismy distribuce odměny prostředí, zatímco tvar pravidel je co nejjednodušší. Prvním pokusem byl tzv. \emph{Kognitivní systém 1 (CS-1)}, který již obsahuje hlavní rysy toho, čemu dnes říkáme Michiganský přístup, protože ho Holland vyvinul na Michiganské univerzitě. Skupina kolem DeJonga a jeho studentů z Pittsburghské univerzity přišla v algoritmu Učící systém 1 (CS-1) s alternativním přístupem, kterému se dnes říká Pittsburghský, a který je orientován na učení s učitelem a klade větší důraz na tvar pravidel a bohatost evolučních operátorů. 

Hlavním rozdílem mezi oběma přístupy je, že v Hollandově případě je jedincem evoluce jedno pravidlo a jako systém pak slouží celá populace pravidel. V Pittsburgském přístupu se naopak uvažuje s jedincem jako kompletním pravidlovým systémem o několika (desítkách) pravidel. Populace pak obsahuje více takových jedinců, jak jsme z evolučních algoritmů zvyklí.  Výsledkem evoluce je nejlepší jedinec, který představuje kompletní řešení problému. 

\subsection{Michiganský přístup: Hollandův LCS}

Jako typický příklad Michiganského přístupu si ukážeme klasický model Hollandova LCS, což je nástupce algoritmu CS-1,který se snaží řešit adaptivní učení agentů, kde zpětná vazba od prostředí přichází jen pro určité stavy a ne po každé akci agenta. 

Jedincem v Hollandově LCS je jedno pravidlo, které má co nejjednodušší tvar. Podmínka je tvořena bitovou maskou pro detekci jednotlivých vstupů/příznaků a může obsahovat hodnoty 0 - příznak není přítomen, 1 - příznak je přítomen a * - na příznaku nezáleží. Vstupy pravidla jsou buď přímo informace z prostředí poskytované receptory, anebo příznaky z interních zpráv, které jsou generovány THEN částí pravidel. Tato pravá strana pravidel tedy obsahuje tzv. zprávu, kterou pravidlo, když je aktivováno, vyšle do globální paměti typu black board. Zpráva slouží jednak k zakódování akcí agenta (např. robot postoupí o krok kupředu) a jednak k přenosu informací mezi pravidly (např. robotu dochází baterie). 

Formálně je jeden jedinec reprezentující pravidlo jednoduchým řetězcem v ternární abecedě, takže genetické operátory křížení a mutace jsou standardní a jednoduché. Holland používá také standardní genetický algoritmus, který ale přechod mezi generacemi realizuje jen zřídka --- po mnoha krocích interakce agentů s prostředím --- a nemění celou populaci, ale nahradí jen nějakou část nejhorších jedinců. 

Pro práci s fitness navrhl Holland komplikovaný algoritmus hasičského družstva (bucket brigade), jenž se snaží řešit problém distribuce odměny posloupnosti akcí, které společně vedou k úspěšné interakci agenta s prostředím. Každý jedinec má kromě pravidla ještě číselnou proměnnou, jejíž hodnota určuje sílu pravidla danou jeho úspěšností.   

Již jsme zmínili, že systém obsahuje sdílenou paměť, která uchovává aktuální zprávy jak od prostředí tak od předchozí aplikace pravidel. V jedné iteraci se zprávy z paměti porovnají se všemi pravidly a vyberou se ta pravidla, jejíchž levá strana odpovídá nějaké zprávě. Tato pravidla vsadí část své síly jako poplatek za účast v soutěži, které pravidlo se skutečně uplatní. Hodnota sázky se počítá jako součin síly pravidla, jeho specificity (počtu 0 a 1 na v podmínce) a konstanty. Pravidlo, které vsadí nejvíce, vyhrává a je aplikováno --- v další iteraci umístí svou zprávu do paměti, ale zároveň musí svou sázku zaplatit. Tato hodnota je pak distribuována mezi pravidla, která byla aplikována v minulých krocích. Tím dochází k tomu, že pravidla, která jsou součástí řetězce akcí vedoucího k dobrému výsledku, posilují (metafora řady hasičů předávající si vědro).

\begin{figure}
\includegraphics{graphics/bucketbrigade.png}
\caption{Hollandův LCS s algoritmem požárního družstva.}
\end{figure}


\subsection{Pittsburghský přístup: GIL}

Jako příklad Pittsburgského přístupu uvedeme sytém GIL, který je praktickým návrhem evoluce binárního klasifikačního systému, který je celý reprezentován jedincem v evolučním algoritmu. Věnujeme se hlavně reprezentaci jedince a operátorům, kterých je velké množství a pracují na různých úrovních.

Jedinec je tedy v tomto systému množina pravidel, které klasifikují do jedné třídy, pravá strana je vždy stejná a proto ji můžeme v jedinci zanedbat. Každé pravidlo si můžeme představit jako disjunkci komplexů, kde každý komplex je konjunkce seletorů z jedné proměnné. Selektor reprezentuje množinu hodnot proměnné a je reprezentován bitovou mapou. 

((X=A1)AND(Z=C3)) OR((X=A2)AND(Y=B2))

[001|11|0011 OR 010|10|1111]

Uveďme stručně výčet operátorů, které se podílejí na úpravách a rekombinacích jedinců, spíše než o mutacích a kříženích mluvíme o úrovních, na kterých operátory pracují.

\begin{description}

\item{Operátory na úrovní jedince:}
výměna pravidel, kopírování pravidel, generalizace pravidla, smazání pravidla, specializace pravidla, zahrnutí jednoho pozitivního příkladu

\item{Operátory na úrovni komplexů:} 
rozdělení komplexu na 1 selektoru, generalizace selektoru (nahrazení 11...1), specializace generalizovaného selektoru, zahrnutí negativního příkladu

\item{Operace na selektorech:}
Mutace 0<->1, rozšíření 0->1, zúžení 1->0, 

\end{description}


\section{ZCS a XCS}

\subsection{ZCS: Klasifikační systém nulté úrovně}

V roce 1994 navrhl Wilson výrazně zjednodušenou verzi LCS, kterou nazval učícím se klasifikátorem nulté úrovně (ZCS). Hlavním rysem algoritmu bylo nahrazení distribuce odměny požárním družstvem pomocí přehledného a jednoduššího algoritmu, a zavedení nového operátoru pokrytí. 

ZCS v každém kroku identifikuje množinu $M$ (matching) pravidel, která jsou aplikovatelná na daný stav prostředí, z nich se vybere to, které se aplikuje čistě na základě fitness jednotlivých pravidel. Všechna pravidla z $M$, která navrhují stejnou akci jako vítězné pravidlo, tvoří pak množinu $A$ (action set). 

Malý zlomek fitness každého pravidla z $A$ je odečten a připraven k redistribuci. Nasčítané části fitness jsou zmenšeny o konstantu GAMMA. Systém si pamatuje členy množiny $A$ z předchozí iterace, mezi které je tato hodnota rovnoměrně rozdělena. 

Pokud systém v aktuálním kroku obdrží odměnu od prostředí, je od ní odečtena malá konstanta a zbytek je rovnoměrně rozdělen mezi členy aktuální množiny $A$. 

Nakonec je ještě zmenšena fitness pravidel z množiny $M-A$ (těch, která jsou aplikovatelná na danou situaci, ale navrhují jinou akci než vítězné pravidlo) o malý faktor $\tau$. 

Vlastní prohledávání je realizováno kombinací stabilního GA (steady-state GA) a operátoru pokrytí. V každém kroku je s určitou pravděpodobností realizována jedna iterace genetického algoritmu, který ruletovou selekcí vybere dva rodiče, vygeneruje dva potomky, a těmi nahradí dva jedince z populace, vybrané opět ruletovou selekcí. Rodiče předají potomkům polovinu své fitness.

Operátor pokrytí se aplikuje tehdy, když je množina M v daném kroku prázdná (nebo pokud obsahuje jen podprůměrné jedince), jeho cílem je zajistit, aby celý systém pokrýval co nejvíce případů, které prostředí poskytuje. Operátor pokrytí vygeneruje pravidlo s podmínkou odpovídající aktuálnímu vstupu a náhodnou akcí na pravé straně. 

\subsection{XCS: Rozšířený klasifikační systém}

Dodnes nejúspěšnějším přístupem v oblasti klasifikačních systémů je druhý algoritmus navržený Wilsonem, tak zvaný rozšířený klasifikační systém (XCS). Autor v něm explicitně aplikoval principy zpětnovazebného Q-učení, které se snaží o úplnost reprezentace problému výběru akce. Algoritmus byl navržen po zkušenostech se ZCS, který, ač byl na řadě úloh úspěšný, vykazuje tendenci soustředit prohledávání jen do oblastí přinášejících velkou odměnu od prostředí. 
Algoritmus XCS nepočítá fitness pravidel na základě hodnoty této odměny, ale na přesnosti předpovědi pravidel. 

Ke každému pravidlu je přiřazena trojice hodnot --- fitness $F$, chyba $\epsilon$ a predikce $p$. Algoritmus zpětnovazebného učení pak upravuje tyto hodnoty následujícím způsobem:  

V každém kroku se identifikuje množina $M$ a pro každou akci $a$ v $M$ se spočte predikce systému jako průměr predikcí pravidel v každé množině $A$ vážený jejich fitness. Odpověď systému se pak stanoví deterministicky nebo s náhodným faktorem. Pokud je množina $M$ prázdná, použije se operátor pokrytí.  

V následujících pěti krocích se pak upraví fitness a další parametry pravidel podle relativní přesnosti pravidla v dané množině $A$.

\begin{enumerate}
\item aktualizuj chybu každého pravidla: $\epsilon_j = \epsilon_j + \beta \left( \left| P - p_j \right| - \epsilon_j \right)$

\item aktualizuj predikci každého pravidla: $p_j = p_j + \beta (P-p_j)$

\item spočti přesnost pravidla $\kappa_j$ jako: $\kappa = \alpha (\epsilon_0/\epsilon)^v$, anebo $\kappa=1$ pro $\epsilon < \epsilon_0$

\item urči relativní přesnot pravidla $\kappa_j\prime$ jako podíl jeho přesnosti a součtu všech přesností pro danou $A$.

\item Aktualizuj fitness $F_j$ na základě relativní přesnosti pomocí moyennovy adaptivní procedury: 

Pokud už fitness byla aktualizována $1/\beta$ krát,
 $F_j = F_j + \beta(\kappa_j\prime - F_j).$ 
 
Jinak nastav $F_j$ jako průměr aktuální a předchozí hodnoty
$kappa_j\prime$.
\end{enumerate}

Nakonec se maximální hodnota $P(a_i)$ zmenší o faktor $\gamma$ a použije se k aktualizaci pravidel z předchozího kroku.

Genetický algoritmus je opět stabilního typu, ale navíc pracuje s faktorizací prostoru pravidel do množin $A$, jde tedy o nikový GA vybírající pravidla z aktuální množiny $A$. Nahrazovaná pravidla jsou vybírána globálně s ohledem na vyvážení velikostí jednotlivých $A$. Pro danou množinu $A$ se rozhoduje o tom, zda se GA aplikuje na základě průměrného času, kdy se naposledy daná nika A podílela na genetickém algoritmu. 

Systém XCS se od původního návrhu dočkal mnoha rozšíření a vylepšení, jako jsou UCS (sUpervised classifier system) vhodný pro učení s učitelem, a jeho další varianty ExSTraCS v. 1 a 2, které byly úspěšné v řešení těžkých praktických problémů. Systém XCSF je zase rozšířením XCS pro problémy aproximace funkcí. 



\chapter{Rojové algoritmy}

Rojové algoritmy mají některé společné rysy s evolučními algoritmy, jako je populační prohledávání, ale jsou charakteistické tím, že kladou důraz na modelování pohybu populace v prostoru parametrů inspirovaném různými druhy organizace sociálního hmyzu nebo jiných živočichů v hejnech. Počáteční výzkum Reynoldse vedl k návrhu jednoduchého lokálního algoritmu tzv. boidů, kteří velice jednoduše a překvapivě dobře modelují organizaci roje a koordinovaný pohyb v něm. Teprve algoritmus optimalizace rojem částic (PSO) ale přinesl využití těchto principů do prohledávacích algoritmů. V současnosti existuje celá řada dalších přibuzných algoritmů od ant colony optimization přes umělé imunitní systémy až ke včelím algoritmům \cite{swarm}. 

Principem fungování rojově optimalizace je organizovaný pohyb roje částic (jedinců) po prostoru parametrů. Každá částice je charakterizována svou polohou a vektorem rychlosti. Každá častice si také pamatuje ze své historie nejlepší místo (ve smyslu hodnoty účelové funkce), které ona sama v historii navštívila a globální dosud nejlepší místo nalezené celým rojem. V každém kroku pak částice upraví svou rychlost jako součet tři vektorů: současné rychlosti, směru k lokálnimu extrému častice a směru ke globálnímu extrému celého roje. Tato suma je vážena dílem vstupními parametry algoritmu nastavovanými experimentátorem (jde o nelehký problém závisející na typu dat) a dílem náhodně vygenreovanými koeficienty pro každý krok. 

\begin{algorithm}
\caption{Schéma optimalizace rojem částic minimalizujících $f:\Real^n\to\Real$}
\label{obrswarm}
\begin{algorithmic}
\Procedure{Optimalizace rojem částic}{}
\State $\vec{g} \gets \arg\max_{\vec{y}} f(\vec{y})$ \Comment inicializace globální nejlepší pozice
\For{$i \gets 1, \dots, l$}
\State Inicializuj pozici částice $\vec{x_i}$ náhodně z rovnoměrného rozdělení
\State $\vec{p_i} \gets \vec{x_i}$ \Comment inicializace nejlepší známé pozice částice
\If{$f(\vec{p_i})<f(\vec{g})$}
\State $\vec{g} \gets \vec{p_i}$ \Comment update globální nejlepší pozice
\EndIf
\State Inicializuj rychlost částice $\vec{v_i}$ náhodně z rovnoměrného rozdělení
\EndFor
\While{$\left(f(\vec{x^t})>\mbox{\emph{kritérium ukončení}}\right)$}
\For{$i \gets 1, \dots, l$}
\State 	zvol $r_g$ a $r_p$ z ronvoměrného rozdělení $U(0,1)$
\For{$d \gets 1, \dots, n$} \Comment update rychlosti částice
\State $v_{i,d} \gets \omega v_{i,d} + \varphi_p r_p (p_{i,d} - x_{i,d}) + \varphi_g r_g (g_{i,d} - x_{i,d})$
\EndFor
\State $\vec{x_i} \gets \vec{x_i} + \vec{v_i}$ \Comment update pozice částice
\If{$f(\vec{x_i})<f(\vec{p_i})$}
\State $\vec{p_i} \gets \vec{x_i}$ \Comment update nejlepší pozice částice
\If{$f(\vec{p_i})<f(\vec{g})$}
\State $\vec{g} \gets \vec{p_i}$ \Comment update nejlepší globální pozice
\EndIf
\EndIf
\EndFor
\EndWhile
\EndProcedure
\end{algorithmic}
\end{algorithm}



\chapter{Genetické programování}
Gentick� programov�n� je zaj�mav� t�m, �e p�in�� pr�ci s jedinci odli�n�ho typu --- jde o stromy reprezentuj�c� programy \cite{koza}. Strom je sestaven z mno�iny netermin�ln�ch symbol� reprezentuj�c�ch funkce a mno�iny termin�ln�ch symbol� (list�) reprezentuj�c�ch prom�nn� a konstanty. Prom�nn� jsou charakterizov�ny �lohou k �e�en�, netermin�ln� symboly ur�uj� programovac� jazyk, ve kter�m budou programy vznikat. Stromy jsou vhodn� jednak proto, �e p�edstavuj� kompaktn� reprezentaci program�, a tak� proto, �e na nich lze elegantn� prov�d�t k���en� realizovan� jako v�m�na n�hodn� zvolen�ch podstrom� mezi dv�ma jedinci. Mutace je realizov�na jako nahrazen� podstromu n�hodn� vygenerovan�m podstromem. Je zaj�mav�, �e v kontrastu s jin�mi oblastmi evolu�n�ch algoritm�, nehraje mutace v genetick�m programov�n� tak v�znamnou roli a obvykle m�v� malou pravd�podobnost. Je to pravd�podobn� zp�sobeno t�m, �e k���en� p�edstavuje ve stromu programu velk� zm�ny, kter� tak supluj� i ��inek mutace.

N�hodn� generov�n� strom� je d�le�itou operac� jak pro inicializaci populace tak pro mutace. Byly navr�eny dv� metody, kter� se sna�� bu� n�hodn� generovat kompletn� strom dan� hloubky, anebo n�hodn� rostou v�tve strom�. Druh� metoda generuje tedy �id�� stromy s nestejn� dlouh�mi v�tvemi. Pro inicializaci se doporu�uje kombinace t�chto metod v pom�ru 1:1 (ramped half-and-half).

\begin{algorithm}
\caption{Sch�ma Kozova algoritmu genetick�ho programov�n� nad syntaktick�mi stromy}
\label{obrsgp}
\begin{algorithmic}
\Procedure{Genetick� programov�n�}{} 
\State $t \gets 0$
\State \emph{Inicializuj} populaci $P_t$ metodou \emph{ramped half-and-half}
\State \emph{Ohodno�} jedince v populaci $P_t$
\While{neplat� \emph{krit�rium ukon�en�}}
\State $i \gets 0$
\Repeat
\State vyber pravd�podobnostn� varia�n� operaci $o\in\{C,M\}$
\If{o=C} \Comment K���en�
\State vyber z $P_t$ 2 rodi�e
\State \emph{Zk�i�} rodi�e
\State vlo� potomky do $P_{t+1}$
\State $i \gets i + 2$
\Else \Comment Mutace
\State vyber z $P_t$ 1 rodi�e
\State \emph{Mutuj} rodi�e
\State vlo� potomka do $P_{t+1}$
\State $i \gets i + 1$
\EndIf
\Until{$i\geq n$}
\State \emph{Ohodno�} jedince v populaci $P_{t+1}$
\State 	Zaho� $P_t$
\State $t \gets t+1$
\EndWhile
%\Until{kr
\EndProcedure
\end{algorithmic}
\end{algorithm}



\chapter{Neuroevoluce}
Použití evolučních algoritmů pro učení neuronových sítí je oblast zkoumaná od přelomu osmdesátých a devadesátých let. První pokusy se týkaly učení vah sítě s předem danou strukturou, jde tedy o využití evolučního algoritmu jako alternativy ke standardnímu učení neuronových sítí založených typicky na lokálních metodách gradientní optimalizace \cite{neuro}. 
Tato úloha není příliš složitá, parametry sítě zakódujeme do reálného vektoru a na něj použijeme jeden z výše popsaných algoritmů. Výpočet fitness pak znamená spočítat chybu neuronové sítě na množině trénovacích dat.

Dnes z experimentů víme, že evoluční algoritmus těžko soupeří v rychlosti s nejlepšími gradientními algoritmy, nicméně jeho výhody pro učení parametrů sítě jsou robustnost proti uváznutí v lokálních extrémech a snadná paralelizace. Jednou z oblastí, kde je evoluční algoritmus pro učení neuronových síti nezastupitelný, jsou případy tzv. posilováného učení, kde není k dispozici informace o chybě sítě pro každý vstup. Tyto případy jsou velmi časté v robotice, kdy teprve z chování robota v delším časovém úseku můžeme odvodit jeho úspěšnost. V této oblasti je evoluční učení neuronových sítí takřka nezastupitelné, i proto se často hovoří o oboru evoluční robotiky. 

Schopnost evoluční algoritmů optimalizovat i složitěji reprezentované struktury vedla ke snahám použít je pro optimalizaci velikosti a propojení jednotek uvnitř sítě, mluvíme o strukturálním učení. Toto učení lze rozdělit podle toho, zda reprezentace sítě je zakódovaná do jedince evolučního algoritmu přímo či nepřímo. Další dělení je, zda se struktura učí najednou společně s parametry sítě (takže kódujeme obě komponenty v jednom genomu) anebo se struktura učí zvlášť a učení parametrů sítě je vlastně záležitostí výpočtu fitness jedince. Při tomto rozdělení úlohy na učení struktury a parametrů je zajímavé si uvědomit, že učit parametry lze libovolným algoritmem, třeba další evolucí.

Při přímých metodách učení struktury sítě se většinou pracuje s reprezentací pomocí matice propojení neuronů v síti. Tato matice může být booleovská (v připadě reprezentace struktury) nebo reálná (pro reprezentaci struktrury i hodnot vah spojů v síti). Jelikož matice propojení může být poměrně velká, vznikly snahy o kompaktní vyjádření struktury sítě, které je často nepřímé. Kitano navrhl gramatické kódování booleovské matice spojů, které uvažuje dvojrozměrnou gramatiku schopnou v logaritmickém prostoru popsat matice navíc s možnosti zachytit symetrie struktury. Gruau navrhl jinou metodu založenou na principu genetického programování, kde program je vlastně návod na sestrojení sítě metodou růstu z minimální konfigurace \cite{gruau}. 

Jednou z dnes nejúspěšnějších neuroevolučních metod je Stanleyho Neat (neuroevolution of augmenting topologies) \cite{stanley}. Jde o přímou metodu vyvíjející zároveň strukturu i parametry sítě opět postupem od minimální konfigurace ke složitějším. Autor elegantně vyřešil problém, jak má vlastně vypadat křížení dvou sítí s různou topologií. Používá k tomu globální paměť evoluční historie tvaru sítě jednoduše realizovatelnou pomocí tzv. rodného čísla spoje, díky kterému lze určit, které hrany v grafu sítě si evolučně odpovídají, a ty se pak mohou křížit. Důležitým problémem při současném učení struktury a parametrů se ukázal krátkodobý negativní vliv strukturálních změn na kvalitu sítě a z toho vyplývající nutnost ochrany strukturálně nových jedinců, kteří potřebují čas na doladění svých parametrů. To je v Neatu opět řešeno pomocí využití rodných čísel hran pro určení podobných sítí a relativizací fitness uvnitř množin podobných sítí. 






\backmatter

\setcounter{tocdepth}{2}
\tableofcontents
\listoffigures
\listoftables
\listofalgorithms 
\printindex
\listoftodos

\nocite{*}
\bibliographystyle{plain}
\bibliography{evabook}

\end{document}
